\section{OpenSSL – részletes parancsok és gyakorlat}

Ez a fejezet az OpenSSL eszközcsomag gyakori parancsait magyarázza el példákkal és biztonsági megjegyzésekkel.

\subsection{Szimmetrikus titkosítás}
\begin{lstlisting}[language=bash]
# AES-128-CBC with password (using \texttt{-k} is not recommended in production)
openssl aes-128-cbc -k PASSWORD -in secret.txt -out secret.txt.aes -p
# Stronger key derivation with PBKDF2
openssl enc -aes-256-cbc -pbkdf2 -iter 200000 -salt -in secret.txt -out secret.txt.aes
\end{lstlisting}

\subsection{RSA kulcsok és konverziók}
\begin{lstlisting}[language=bash]
# Generate RSA 2048 key pair
openssl genpkey -algorithm RSA -out private.pem -pkeyopt rsa_keygen_bits:2048
# Extract public key
openssl rsa -pubout -in private.pem -out public.pem
# PKCS\#8 conversion (unencrypted)
openssl pkcs8 -topk8 -inform PEM -outform PEM -nocrypt -in private.pem -out private-pkcs8.pem
\end{lstlisting}

\subsection{Aláírás és ellenőrzés}
\begin{lstlisting}[language=bash]
# Create a SHA256 signature
openssl dgst -sha256 -sign private.pem -out doc.sign doc.txt
# Verify with the public key
openssl dgst -sha256 -verify public.pem -signature doc.sign doc.txt
\end{lstlisting}

\subsection{Titkosítás OAEP-paddal}
\begin{lstlisting}[language=bash]
openssl pkeyutl -encrypt -inkey public.pem -pubin -in secret.txt -out secret.txt.enc -pkeyopt rsa_padding_mode:oaep -pkeyopt rsa_oaep_md:sha256
openssl pkeyutl -decrypt -inkey private.pem -in secret.txt.enc -out secret.txt.dec
\end{lstlisting}

Megjegyzés: mindig használjunk modern padding-ot (OAEP) és elegendő kulcshosszt (>=2048). A PKCS\#8 formátum modern alapértelmezés a kulcsok cseréjéhez.
